\begin{prob}{009}{}{2023年1月号 学コン6番}
 表と裏が等確率で出るコインを$n$回投げて, 点${\rm P}_n(x_n,y_n,z_n)$を次のように定める.
 \begin{align*}
 &(x_0,y_0,z_0) = (1,-3,1) \\
 &\text{$k$回目で表が出たとき}(x_k,y_k,z_k) = (-x_{k-1}, -y_{k-1} - z_{k-1} , x_{k-1}+y_{k-1} + z_{k-1}) \\
 &\text{$k$回目で裏が出たとき}(x_k,y_k,z_k) = (z_{k-1}, y_{k-1}, x_{k-1})
 \end{align*}
 \begin{enumerate}
  \item ${\rm P}_n$が$(1,-1,0)$である確率を$a_n$, $(-1,1,0)$である確率を$b_n$として$a_n+b_n$を$n$で表せ.
  \item ${\rm P}_n$が$(0,0,0)$である確率を$n$で表せ.
 \end{enumerate}
\end{prob}
